\documentclass[11pt]{article}

\usepackage[margin=1in]{geometry}
\usepackage{amsmath,amssymb,amsthm,mathtools,bm}
\usepackage{microtype}
\usepackage{enumitem}
\usepackage[hidelinks]{hyperref}

\newtheorem{theorem}{Theorem}[section]
\newtheorem{proposition}[theorem]{Proposition}
\newtheorem{lemma}[theorem]{Lemma}
\newtheorem{corollary}[theorem]{Corollary}
\theoremstyle{remark}
\newtheorem{remark}[theorem]{Remark}

\newcommand{\R}{\mathbb R}
\newcommand{\E}{\mathbb E}
\newcommand{\Pp}{\mathbb P}
\newcommand{\ip}[2]{\left\langle #1,#2\right\rangle}
\newcommand{\norm}[1]{\left\lVert #1\right\rVert}
\newcommand{\op}{\mathrm{op}}
\newcommand{\F}{\mathrm{F}}
\newcommand{\1}{\mathbf 1}
\newcommand{\diag}{\operatorname{diag}}
\newcommand{\tr}{\operatorname{tr}}
\newcommand{\vecop}{\operatorname{vec}}
\newcommand{\psiOne}{\psi_1}
\newcommand{\psiTwo}{\psi_2}

\title{Sharp Frobenius Radius of the Gaussian Softmax\\
Associative-Memory Minimizer for Zipf Frequencies\\
$0\leq \alpha\leq 1$ Above Capacity}
\author{Technical note}
\date{July 2026}

\begin{document}
\maketitle

\begin{abstract}
Let $N=\lceil d^\beta\rceil$ with fixed $\beta>2$, and let
$p_i=i^{-\alpha}/H_{N,\alpha}$ with fixed $0\leq\alpha\leq1$.
For the Gaussian softmax associative-memory loss, let $W_\star$ denote
its unique finite minimizer on the event supplied by the previously
proved existence-and-uniqueness theorem.  We prove that, with high
probability,
\[
 \norm{W_\star}_{\F}
 =\widetilde\Theta_{\alpha,\beta}\!\left(
 d^2\Bigg[\sum_{i=1}^N\bigl(p_i\wedge d^{-2}\bigr)^2\Bigg]^{1/2}
 \right).
\]
The proof is uniform in the adaptive minimizer.  For $\alpha<1$, the
upper bound follows from coercivity on a deterministic diffuse tail,
while the lower bound follows from a deterministic clipped comparison
matrix and a deterministic weighted capped-support inequality.  The
endpoint $\alpha=1$ is treated by a head-memory comparison and a
blockwise coercivity estimate.  All nonstandard probabilistic claims
used in these arguments are proved below; only standard concentration
inequalities and the previously established existence-and-uniqueness
theorem are imported.
\end{abstract}

\section{Model and main result}

Fix constants
\[
 0\leq\alpha\leq1,
 \qquad \beta>2,
 \qquad N=\lceil d^\beta\rceil,
 \qquad D=d^2,
 \qquad h=\log N.
\]
Let
\[
 u_1,\ldots,u_N,v_1,\ldots,v_N
 \overset{\mathrm{ind}}{\sim}\mathcal N(0,I_d/d),
 \qquad
 p_i=\frac{i^{-\alpha}}{H_{N,\alpha}},
 \qquad
 H_{N,\alpha}=\sum_{k=1}^N k^{-\alpha}.
\]
For $W\in\R^{d\times d}$, define
\[
 \ell_i(W)
 =\log\!\left(1+\sum_{j\neq i}
 \exp\bigl((u_j-u_i)^\top Wv_i\bigr)\right),
 \qquad
 L(W)=\sum_{i=1}^N p_i\ell_i(W).
\]
We use the following theorem proved in the companion existence note.

\begin{theorem}[Imported existence and uniqueness theorem]
\label{thm:existence}
If $N>D$, define
\[
 \mathcal E_{\rm ex}
 :=\{L\text{ has a unique finite minimizer in }\R^{d\times d}\}.
\]
Then
\[
 \Pp(\mathcal E_{\rm ex})
 \geq
 1-2^{1-N}\sum_{k=0}^{D-1}\binom{N-1}{k}.
\]
For $N=\lceil d^\beta\rceil$ with $\beta>2$, there exists
$c_{\rm ex}=c_{\rm ex}(\beta)>0$ such that, for all sufficiently large
$d$,
\[
 \Pp(\mathcal E_{\rm ex}^c)\leq e^{-c_{\rm ex}N}.
\]
\end{theorem}

For an outcome $\omega$, write $L_\omega$ for the corresponding
random loss and define
\[
 W_\star(\omega)
 :=
 \begin{cases}
 \displaystyle\arg\min_{W\in\R^{d\times d}}L_\omega(W),
 &\omega\in\mathcal E_{\rm ex},\\[2mm]
 0,&\omega\notin\mathcal E_{\rm ex}.
 \end{cases}
\]
This is a measurable random matrix: $L_\omega(W)$ is jointly measurable
in $\omega$, continuous in $W$, and on $\mathcal E_{\rm ex}$ has a
unique attained minimum, so the measurable-argmin theorem applies
(after localization to increasing compact balls).  Every use of
optimality below is explicitly restricted to $\mathcal E_{\rm ex}$;
the arbitrary extension on its complement has no mathematical role.

Define the clipped Frobenius scale
\begin{equation}
 R_F
 :=D\left(\sum_{i=1}^N (p_i\wedge D^{-1})^2\right)^{1/2}.
 \label{eq:RF-def}
\end{equation}

\begin{theorem}[Sharp Frobenius scale]
\label{thm:main}
For every fixed $0\leq\alpha\leq1$ and $\beta>2$, there exist constants
$c,C,c_0>0$, depending only on $(\alpha,\beta)$, such that, for all
sufficiently large $d$,
\[
 \Pp\!\left(
 \mathcal E_{\rm ex}\cap
 \left\{h^{-C}R_F\leq\norm{W_\star}_{\F}\leq h^C R_F\right\}
 \right)
 \geq 1-d^{-c_0}-e^{-cN}.
\]
Consequently,
\[
 \norm{W_\star}_{\F}=\widetilde\Theta_{\alpha,\beta}(R_F).
\]
More explicitly,
\[
 R_F\asymp_{\alpha,\beta}
 \begin{cases}
 d^{2-\beta/2},
   &0\leq\alpha<\frac12,\\[2mm]
 d^{2-\beta/2}\sqrt{\log N},
   &\alpha=\frac12,\\[2mm]
 d^{2-\beta(1-\alpha)},
   &\frac12<\alpha<1,\ \beta(1-\alpha)\geq2,\\[2mm]
 d^{\frac{2-\beta(1-\alpha)}{2\alpha}},
   &\frac12<\alpha<1,\ \beta(1-\alpha)<2,\\[2mm]
 d/\sqrt{\log N},
   &\alpha=1.
 \end{cases}
\]
\end{theorem}

All subsequent assertions are understood for sufficiently large $d$.
Constants denoted by $c,C$ may change from line to line; dependence on
fixed $(\alpha,\beta)$ is suppressed.

\section{Centered representation and deterministic inequalities}

Set
\[
 \bar u=\frac1N\sum_{j=1}^N u_j,
 \qquad
 x_j=u_j-\bar u,
 \qquad
 \Phi(y)=\log\!\left(\frac1N\sum_{j=1}^N e^{x_j^\top y}\right).
\]
Since $\sum_jx_j=0$, Jensen's inequality gives $\Phi(y)\geq0$.  Moreover,
\begin{equation}
 \ell_i(W)=h+\Phi(Wv_i)-x_i^\top Wv_i.
 \label{eq:centered-loss}
\end{equation}
Indeed,
\[
 1+\sum_{j\neq i}e^{(u_j-u_i)^\top Wv_i}
 =\sum_{j=1}^Ne^{(x_j-x_i)^\top Wv_i}
 =N e^{-x_i^\top Wv_i+\Phi(Wv_i)}.
\]
For a deterministic set $T\subseteq[N]$, write
\[
 P_T=\sum_{i\in T}p_i,
 \qquad
 S_T=\sum_{i\in T}p_i^2,
 \qquad
 G_T=\sum_{i\in T}p_i x_iv_i^\top.
\]
Because every loss is nonnegative, summing only over $T$ in
\eqref{eq:centered-loss} yields
\begin{equation}
 L(W)\geq P_T h+\sum_{i\in T}p_i\Phi(Wv_i)-\ip{G_T}{W}_{\F}.
 \label{eq:subset-lower}
\end{equation}
Define the improvement over $W=0$ by
\[
 I(W)=h-L(W)
 =\sum_{i=1}^Np_i\bigl(x_i^\top Wv_i-\Phi(Wv_i)\bigr).
\]
Since $\Phi\geq0$ and $\ell_i(W)\geq0$, if
$s_i(W)=x_i^\top Wv_i$, then
\begin{equation}
 p_i\bigl(s_i(W)-\Phi(Wv_i)\bigr)
 \leq p_i\min\{h,(s_i(W))_+\}.
 \label{eq:deterministic-cap}
\end{equation}
This deterministic cap is the only clipping statement used for the
adaptive minimizer.

\section{Zipf arithmetic for $0\leq\alpha<1$}

Throughout this section, $0\leq\alpha<1$.  Integral comparison gives
\begin{equation}
 H_{N,\alpha}\asymp_\alpha N^{1-\alpha},
 \qquad
 p_i\asymp_\alpha N^{\alpha-1}i^{-\alpha}.
 \label{eq:zipf-basic}
\end{equation}
Set
\begin{equation}
 b_i=\min\{1,Dp_i\},
 \qquad
 R^2=\sum_{i=1}^N b_i^2,
 \qquad
 J=\sum_{i=1}^Np_ib_i.
 \label{eq:bRJ}
\end{equation}
By definition, $R=R_F$.  Let
\[
 m=\max\{i:p_i\geq D^{-1}\},
\]
with $m=0$ if the set is empty.  Then
\begin{equation}
 R^2=m+D^2\sum_{i>m}p_i^2.
 \label{eq:R-decomp}
\end{equation}

\begin{lemma}[Clipped energy]
\label{lem:arithmetic}
For fixed $0\leq\alpha<1$ and $\beta>2$,
\begin{align}
 J&\asymp_\alpha R^2/D,
 \label{eq:J-R}\\
 \sum_{i\leq m}p_i&\asymp_\alpha m/D
 \quad\text{when }m\geq1,
 \label{eq:head-mass}\\
 \frac{\norm{p}_2R}{dJ}&\leq d^{-c}h^C
 \quad\text{for some }c>0,
 \label{eq:cross-small}\\
 R/d&\leq Ch^{1/2}d^{-\eta}
 \quad\text{for some }\eta=\eta(\alpha,\beta)>0.
 \label{eq:R-small}
\end{align}
Consequently, for every fixed $B>0$,
\begin{equation}
 h^B R/d\longrightarrow0.
 \label{eq:R-polylog-small}
\end{equation}
\end{lemma}

\begin{proof}
For $i>m$, $b_i=Dp_i$, hence
\[
 \sum_{i>m}p_ib_i
 =D\sum_{i>m}p_i^2
 =\frac1D\sum_{i>m}b_i^2.
\]
If $m=0$, this proves \eqref{eq:J-R}.  Suppose $m\geq1$.  Notice first
that $m<N$: indeed,
$p_N=N^{-\alpha}/H_{N,\alpha}\leq N^{-1}<D^{-1}$ because $N>D$.
Monotonicity and $p_{m+1}<D^{-1}\leq p_m$ imply
$D^{-1}\leq p_m\leq 2^\alpha D^{-1}$.  Therefore
\[
 \sum_{i\leq m}p_i\asymp_\alpha mp_m\asymp_\alpha m/D,
\]
which proves \eqref{eq:head-mass}; comparing the head and tail
contributions proves \eqref{eq:J-R}.

If $m=0$, then $R=D\norm p_2$ and $J=D\norm p_2^2$, so
\[
 \frac{\norm p_2R}{dJ}=\frac1d.
\]
Now assume $m\geq1$.

If $0\leq\alpha<1/2$, then $\norm p_2\asymp N^{-1/2}$.  Moreover,
$p_m\asymp D^{-1}$ gives
\[
 m\asymp(DN^{\alpha-1})^{1/\alpha}
 \quad(\alpha>0),
\]
and
\[
 \frac{m}{D^2/N}
 \asymp d^{\frac{(\beta-2)(2\alpha-1)}{\alpha}}=o(1).
\]
Thus the tail in \eqref{eq:R-decomp} satisfies
$D^2\sum_{i>m}p_i^2\asymp D^2/N$, so
$R\asymp D/\sqrt N$ and $J\asymp D/N$.  Hence the ratio in
\eqref{eq:cross-small} is $O(d^{-1})$.  The case $\alpha=0$ cannot
have $m\geq1$ for large $d$ because $N>D$.

If $\alpha=1/2$, then
$\norm p_2\asymp\sqrt{h/N}$ and, whether or not clipping occurs,
\[
 R\asymp D\sqrt{h/N},
 \qquad
 J\asymp Dh/N.
\]
Indeed, in the clipped case $m\asymp D^2/N$ and
$\sum_{i>m}p_i^2\asymp h/N$.  Again
\eqref{eq:cross-small} is $O(d^{-1})$ up to an absolute logarithmic
factor.

Finally let $1/2<\alpha<1$.  Since
$\sum_i i^{-2\alpha}\asymp1$, one has
$\norm p_2\asymp p_1$.  The relation $p_m\asymp D^{-1}$ gives
$p_1\asymp m^\alpha/D$.  Also
\[
 D^2\sum_{i>m}p_i^2\asymp m,
 \qquad R\asymp\sqrt m,
 \qquad J\asymp m/D.
\]
Consequently,
\[
 \frac{\norm p_2R}{dJ}\asymp\frac{m^{\alpha-1/2}}d.
\]
If $m$ grows polynomially, write
\[
 m\asymp d^\kappa,
 \qquad
 \kappa=\frac{2-\beta(1-\alpha)}{\alpha}<2.
\]
Then
\[
 \kappa\left(\alpha-\frac12\right)-1
 =\frac{(\alpha-1)(2\alpha\beta-\beta+2)}{2\alpha}<0.
\]
If $m$ is bounded, the ratio is $O(d^{-1})$.  This proves
\eqref{eq:cross-small}.

It remains to prove \eqref{eq:R-small}.  If $\alpha\leq1/2$, the
preceding estimates give
\[
 \frac Rd\lesssim d^{1-\beta/2}\sqrt h,
\]
which remains $o(1)$ after multiplication by any fixed power of $h$
because $\beta>2$.  If $1/2<\alpha<1$ and $m=0$, then
$R=D\norm{p}_2=O(1)$, so $R/d=O(d^{-1})$.  If $m\geq1$ is bounded,
the same conclusion holds.  Finally, if $m$ grows polynomially, then
$R\asymp d^{\kappa/2}$ and
\[
 1-\frac\kappa2
 =\frac{(\beta-2)(1-\alpha)}{2\alpha}>0.
\]
Thus, in every case, there is an
$\eta=\eta(\alpha,\beta)>0$ such that
\[
 \frac Rd\leq Ch^{1/2}d^{-\eta},
\]
which proves \eqref{eq:R-small}.  Since $h=\beta\log d+o(1)$, every
fixed power of $h$ is dominated by $d^{\eta/2}$ for all sufficiently
large $d$, and \eqref{eq:R-polylog-small} follows.
\end{proof}

We also need a tail on which every normalized weight is much smaller
than $D^{-1}$.  Fix once and for all
\[
 A=24,
 \qquad
 T=\{i:p_i\leq(Dh^A)^{-1}\},
 \qquad
 H=[N]\setminus T.
\]

\begin{lemma}[Strongly diffuse tail]
\label{lem:strong-tail}
For fixed $0\leq\alpha<1$ and $\beta>2$,
\begin{align}
 P_T&\geq\frac12,
 &P_Hh&=o(1),
 &d\sqrt{S_T}&\leq h^{-A/2},
 \label{eq:strong-tail-basic}\\
 D\sqrt{S_T}+d\sqrt{P_Hh}&\leq h^C R
 \label{eq:strong-tail-compare}
\end{align}
for sufficiently large $d$ and a constant $C=C(\alpha,\beta)$.
\end{lemma}

\begin{proof}
For $\alpha=0$, $p_i=N^{-1}<(Dh^A)^{-1}$ for large $d$, so $H$ is
empty.  Assume $0<\alpha<1$ and let $r=|H|$.  If $r=0$, then
$H=\varnothing$, $P_T=1$, $P_H=0$, and
\[
 S_T\leq(Dh^A)^{-1},
 \qquad
 D^2S_T=\sum_i(Dp_i)^2\leq R^2;
\]
hence all the asserted inequalities are immediate.  We may therefore
assume $r\geq1$.  Then \eqref{eq:zipf-basic} and the definition of $H$
give
\begin{equation}
 r\asymp_\alpha(Dh^A N^{\alpha-1})^{1/\alpha},
 \qquad
 \frac rD\lesssim
 h^{A/\alpha}d^{-\frac{(\beta-2)(1-\alpha)}{\alpha}}=o(1).
 \label{eq:strong-head-size}
\end{equation}
Regular variation gives
\[
 P_H\asymp_\alpha r p_r\asymp_\alpha\frac{r}{Dh^A}.
\]
Together with \eqref{eq:strong-head-size}, this implies $P_Hh=o(1)$
and hence $P_T\geq1/2$.  Since $p_i\leq(Dh^A)^{-1}$ on $T$,
\[
 S_T\leq \max_{i\in T}p_i\,P_T\leq(Dh^A)^{-1},
\]
which proves the third claim in \eqref{eq:strong-tail-basic}.
Moreover, $p_i<D^{-1}$ on $T$, so
\[
 D^2S_T=\sum_{i\in T}(Dp_i)^2\leq R^2.
\]
Thus the first term in \eqref{eq:strong-tail-compare} is at most $R$.

It remains to control $d\sqrt{P_Hh}$.  If the natural clipping index
$m$ is positive, then comparison of the thresholds $D^{-1}$ and
$(Dh^A)^{-1}$ gives $r\leq C h^{A/\alpha}m$.  Hence
\[
 d\sqrt{P_Hh}
 \lesssim \sqrt r\,h^{(1-A)/2}
 \leq h^{A/(2\alpha)+(1-A)/2}\sqrt m
 \leq h^C R.
\]
Suppose $m=0$.  If $2-\beta(1-\alpha)<0$, then
$Dp_1\leq Cd^{2-\beta(1-\alpha)}<h^{-A}$ for large $d$, so $H$ is
empty.  If $2-\beta(1-\alpha)=0$, then $r\leq Ch^{A/\alpha}$ and
$d\sqrt{P_Hh}\leq h^C$.  In this critical case, $R=D\norm p_2$ is
bounded below by a positive constant when $\alpha>1/2$, is of order
$\sqrt h$ when $\alpha=1/2$, and grows polynomially when
$\alpha<1/2$.  Therefore $d\sqrt{P_Hh}\leq h^C R$.  The remaining
case $2-\beta(1-\alpha)>0$ forces $m\geq1$ for large $d$ and was
already treated.
\end{proof}

\section{Concentration tools}

For $r\in\{1,2\}$, we use the Orlicz norms
\[
 \norm X_{\psi_r}
 :=\inf\left\{s>0:
 \E\exp\!\left(\frac{|X|^r}{s^r}\right)\leq2\right\}.
\]
Conditional Orlicz norms are defined by replacing $\E$ with conditional
expectation; they are measurable with respect to the conditioning
sigma-field.

We use the following standard inequalities in their usual forms.
They are stated to make every later parameter choice explicit.

\begin{lemma}[Standard concentration inequalities]
\label{lem:standard-tools}
There are absolute constants $c,C>0$ such that the following hold.
\begin{enumerate}[label=(\roman*)]
\item If $X$ is a centered polynomial of total degree at most $q_0$ in
independent standard Gaussian variables, then, for every $q\geq2$,
\[
 \norm X_{L^q}\leq(q-1)^{q_0/2}\norm X_{L^2}.
\]
\item If $X_1,\ldots,X_n$ are independent centered subexponential
random variables, put
\[
 K_i=\norm{X_i}_{\psiOne},\qquad
 V^2=\sum_iK_i^2,\qquad K_\ast=\max_iK_i.
\]
Then, for every $t\geq1$,
\[
 \Pp\!\left(\left|\sum_iX_i\right|
 >C(V\sqrt t+K_\ast t)\right)\leq2e^{-t}.
\]
The same assertion applies under a conditional law whenever the
variables are conditionally independent and centered and the displayed
conditional Orlicz-norm bounds hold almost surely.
\item If $Y_i$ are independent centered self-adjoint $M\times M$
matrices with $\norm{Y_i}_{\op}\leq L$ almost surely and
$\sigma^2=\norm{\sum_i\E Y_i^2}_{\op}$, then
\[
 \Pp\!\left(\norm{\sum_iY_i}_{\op}\geq t\right)
 \leq2M\exp\!\left(-\frac{t^2/2}{\sigma^2+Lt/3}\right).
\]
\item If $g\sim\mathcal N(0,I_m)$ and $B\succeq0$, then, for every
$x>0$,
\[
 \Pp\bigl(g^\top Bg-\tr B
 \geq2\norm B_{\F}\sqrt x+2\norm B_{\op}x\bigr)\leq e^{-x}.
\]
\item If $F:\R^M\to\R$ is $L$-Lipschitz and $g\sim\mathcal N(0,I_M)$,
then
\[
 \Pp(F(g)\geq\E F(g)+t)\leq e^{-t^2/(2L^2)}.
\]
\item If $G$ is a $d\times n$ standard Gaussian matrix, then, for
$t\geq0$,
\[
 \Pp\bigl(\norm G_{\op}>\sqrt n+\sqrt d+t\bigr)\leq e^{-t^2/2}.
\]
\end{enumerate}
\end{lemma}

We repeatedly use the elementary Gaussian tail consequence
\begin{equation}
 \E\bigl[(1+\norm v^8)\1_{\{\norm v>2\}}\bigr]\leq Ce^{-cd},
 \qquad v\sim\mathcal N(0,I_d/d),
 \label{eq:gaussian-trunc-tail}
\end{equation}
and the fact that, for $q\leq cd$ and independent
$v,w\sim\mathcal N(0,I_d/d)$,
\begin{equation}
 \norm{v^\top w}_{L^q}\leq C\sqrt{q/d}.
 \label{eq:inner-product-moment}
\end{equation}
Both follow directly from the chi-square tail bound and conditioning.
We also use the elementary Orlicz-product inequality
\begin{equation}
 \norm{XY}_{\psiOne}
 \leq C\norm X_{\psiTwo}\norm Y_{\psiTwo}.
 \label{eq:orlicz-product}
\end{equation}
Indeed, after normalizing the two $\psi_2$ norms, this follows from
$2|XY|\leq X^2+Y^2$ and Cauchy--Schwarz.  The same proof applies under
a conditional law.

\subsection{Uniform lower bound for the empirical key free energy}

Let
\[
 \psi(t)=\min\{t^2,t\},\qquad t\geq0.
\]

\begin{lemma}[Uniform empirical free-energy lower bound]
\label{lem:free-energy}
With probability at least $1-e^{-cN}-e^{-cd}$,
\begin{equation}
 \Phi(y)\geq c\psi\!\left(\frac{\norm y_2}{\sqrt d}\right)
 \qquad\text{for every }y\in\R^d.
 \label{eq:free-energy-lower}
\end{equation}
\end{lemma}

\begin{proof}
Choose a nonzero nonnegative Lipschitz function
$\rho:\R\to[0,1]$, supported on $[-3,3]$, such that
$\rho(z)\leq z^2$ for every $z$.  Put
$c_\rho=\E\rho(g)>0$ for $g\sim\mathcal N(0,1)$, and let
$L_\rho$ be the Lipschitz constant.  Let $\mathcal N$ be a
$\delta$-net of the unit sphere with
$\delta=c_\rho/(64L_\rho\sqrt d)$.  Then
$|\mathcal N|\leq(C\sqrt d)^d$.

For fixed $a\in\mathcal N$, the variables
$\rho(\sqrt d\,a^\top u_j)$ are i.i.d., bounded by one, and have mean
$c_\rho$.  Hoeffding's inequality gives
\[
 \Pp\!\left(\frac1N\sum_{j=1}^N
 \rho(\sqrt d\,a^\top u_j)<\frac{3c_\rho}{4}\right)
 \leq e^{-cN}.
\]
Because $N=d^\beta$ with $\beta>2$, the union bound over
$\mathcal N$ is still at most $e^{-cN}$.  Also,
\[
 \Pp\!\left(\max_j\norm{u_j}_2>2\right)\leq e^{-cd},
 \qquad
 \Pp\!\left(\sqrt d\,\norm{\bar u}_2>\frac{c_\rho}{32L_\rho}
 \right)\leq e^{-cN}.
\]
On the intersection of these events, for every unit vector $a$, choose
$a_0\in\mathcal N$ with $\norm{a-a_0}_2\leq\delta$.  Then, for every
$j$,
\[
 \left|\sqrt d\,a^\top x_j-
 \sqrt d\,a_0^\top u_j\right|
 \leq\sqrt d\,\norm{a-a_0}_2\norm{u_j}_2
      +\sqrt d\,\norm{\bar u}_2
 \leq\frac{c_\rho}{16L_\rho}.
\]
The Lipschitz property therefore implies
\begin{equation}
 \frac1N\sum_{j=1}^N\rho(\sqrt d\,a^\top x_j)
 \geq\frac{c_\rho}{2}
 \qquad\text{for every unit }a.
 \label{eq:truncated-variance}
\end{equation}

Fix a unit vector $a$ and set
$f_a(t)=\Phi(t\sqrt d\,a)$ for $t\geq0$.  Since $\sum_jx_j=0$,
$f_a(0)=f_a'(0)=0$.  There is an absolute $c_1>0$ such that
\[
 e^z\geq1+z+c_1z^2\1_{\{|z|\leq1\}}
 \qquad(z\in\R).
\]
For $0\leq t\leq1/3$, apply this inequality to
$z_j=t\sqrt d\,a^\top x_j$.  Since $\rho$ is supported on $[-3,3]$,
\eqref{eq:truncated-variance} gives
\[
 e^{f_a(t)}
 =\frac1N\sum_je^{z_j}
 \geq1+c_1t^2\frac1N\sum_j\rho(\sqrt d\,a^\top x_j)
 \geq1+ct^2.
\]
Hence $f_a(t)\geq ct^2$ for $0\leq t\leq1/3$.  Since $f_a$ is convex
and $f_a(0)=0$, $f_a(t)/t$ is nondecreasing on $(0,\infty)$; therefore
$f_a(t)\geq ct$ for $t\geq1/3$.  Writing
$y=t\sqrt d\,a$ proves \eqref{eq:free-energy-lower}.
\end{proof}

\subsection{Uniform weighted query regularity}

\begin{lemma}[Diffuse weighted query moments]
\label{lem:weighted-query}
Let $q_i\geq0$ be deterministic, $\sum_iq_i=1$, and
\begin{equation}
 q_{\max}:=\max_iq_i\leq(Dh^{12})^{-1}.
 \label{eq:qmax}
\end{equation}
Then, with probability at least $1-d^{-30}$,
\begin{align}
 \frac1{2d}I_d
 &\preceq\sum_iq_iv_iv_i^\top
 \preceq\frac2dI_d,
 \label{eq:weighted-cov}\\
 \sum_iq_i(v_i^\top Av_i)^2
 &\leq\frac{C}{d^2}(\tr A)^2
 \qquad\text{for every }A\succeq0.
 \label{eq:weighted-fourth}
\end{align}
Consequently, on the same event,
\begin{equation}
 \sum_iq_i\psi\!\left(\frac{\norm{Wv_i}_2}{\sqrt d}\right)
 \geq c\psi\!\left(\frac{\norm W_{\F}}d\right)
 \qquad\text{for every }W\in\R^{d\times d}.
 \label{eq:weighted-coercivity}
\end{equation}
\end{lemma}

\begin{proof}
We first prove \eqref{eq:weighted-cov}.  Put
$Y_i=q_i(v_iv_i^\top-I_d/d)$ and
$E_i=\{\norm{v_i}_2\leq2\}$.  Apply matrix Bernstein to the centered,
bounded matrices
\[
 \widehat Y_i=Y_i\1_{E_i}-\E[Y_i\1_{E_i}].
\]
On $E_i$, $\norm{\widehat Y_i}_{\op}\leq Cq_{\max}$.  Moreover,
\[
 \left\|\sum_i\E\widehat Y_i^2\right\|_{\op}
 \leq C\sum_iq_i^2
 \left\|\E(v_iv_i^\top-I_d/d)^2\right\|_{\op}
 \leq Cq_{\max}/d.
\]
The truncation bias satisfies, by \eqref{eq:gaussian-trunc-tail},
\[
 \left\|\sum_i\E[Y_i\1_{E_i^c}]\right\|_{\op}
 \leq Ce^{-cd}.
\]
Taking matrix-Bernstein deviation $1/(4d)$ gives exponent
$c/(dq_{\max})\geq cdh^{12}$.  Together with
$\Pp(\cup_iE_i^c)\leq e^{-cd}$, this proves
\eqref{eq:weighted-cov} with failure much smaller than $d^{-30}$.

For the fourth moment, work in the Hilbert space
$\mathcal H$ of symmetric $d\times d$ matrices, of dimension at most
$D$.  Put
\[
 Z_i=v_iv_i^\top-I_d/d,
 \qquad
 \mathcal A_i=Z_i\otimes Z_i,
\]
where $(Z\otimes Z)A=\ip{Z}{A}_{\F}Z$.  For symmetric $A$,
Gaussian fourth moments give
\begin{equation}
 \E\ip{Z_i}{A}_{\F}^2=\frac2{d^2}\norm A_{\F}^2.
 \label{eq:Z-cov}
\end{equation}
Furthermore,
\begin{align*}
 &\ip{A}{\E[\norm{Z_i}_{\F}^2(Z_i\otimes Z_i)]A}_{\F}\\
 &\qquad=\E\bigl[\norm{Z_i}_{\F}^2\ip{Z_i}{A}_{\F}^2\bigr]\\
 &\qquad\leq
 (\E\norm{Z_i}_{\F}^4)^{1/2}
 (\E\ip{Z_i}{A}_{\F}^4)^{1/2}
 \leq\frac{C}{d^2}\norm A_{\F}^2.
\end{align*}
The last inequality uses the degree-two Gaussian hypercontractive
bound together with \eqref{eq:Z-cov}.  Hence
\begin{equation}
 \left\|\E[\norm{Z_i}_{\F}^2(Z_i\otimes Z_i)]\right\|_{\op}
 \leq C/d^2.
 \label{eq:operator-fourth-variance}
\end{equation}

Truncate again on $E_i$.  On $E_i$,
$\norm{\mathcal A_i}_{\op}=\norm{Z_i}_{\F}^2\leq C$.  By
\eqref{eq:operator-fourth-variance}, the variance proxy for
\[
 q_i\bigl(\mathcal A_i\1_{E_i}-
 \E[\mathcal A_i\1_{E_i}]\bigr)
\]
is at most $Cq_{\max}/d^2$, while the summand norm is at most
$Cq_{\max}$.  Matrix Bernstein with deviation $d^{-2}$ has exponent
$c/(Dq_{\max})\geq ch^{12}$.  The dimension factor is at most $2D$, and the truncation bias is
$O(e^{-cd})$ by \eqref{eq:gaussian-trunc-tail}.  Moreover,
\[
 \Pp\!\left(\bigcup_iE_i^c\right)
 \leq N\Pp(\norm{v_1}_2>2)\leq e^{-c'd},
\]
because $N=d^\beta+O(1)$.  On $\cap_iE_i$, the truncated and
untruncated empirical operators coincide.  Combining the operator
deviation, its mean, the truncation bias, and this last event shows
that, with failure at most $d^{-30}$,
\[
 \sum_iq_i\ip{Z_i}{A}_{\F}^2
 \leq\frac{C}{d^2}\norm A_{\F}^2
 \qquad\text{for every symmetric }A.
\]
For $A\succeq0$,
\[
 v_i^\top Av_i=\frac{\tr A}{d}+\ip{Z_i}{A}_{\F},
\]
so $(a+b)^2\leq2a^2+2b^2$ and
$\norm A_{\F}\leq\tr A$ prove \eqref{eq:weighted-fourth}.

Finally, fix $W$ and set $r=\norm W_{\F}$.  If $r=0$, then
\eqref{eq:weighted-coercivity} is immediate.  Assume $r>0$.  If an
index $I$ is drawn with distribution $(q_i)$, then \eqref{eq:weighted-cov} and
\eqref{eq:weighted-fourth}, applied to $A=W^\top W$, give
\[
 \E_q\norm{Wv_I}_2^2\geq\frac{r^2}{2d},
 \qquad
 \E_q\norm{Wv_I}_2^4\leq C\frac{r^4}{d^2}.
\]
Paley--Zygmund therefore gives a fixed $c_0>0$ such that a $q$-mass at
least $c_0$ satisfies
$\norm{Wv_i}_2^2\geq r^2/(4d)$.  On this set,
\[
 \psi\!\left(\frac{\norm{Wv_i}_2}{\sqrt d}\right)
 \geq c\psi(r/d).
\]
Averaging proves \eqref{eq:weighted-coercivity}.
\end{proof}

\begin{corollary}[Diffuse-tail coercivity]
\label{cor:tail-coercivity}
Let $T\subseteq[N]$ be deterministic, $P_T\geq1/2$, and
\[
 \max_{i\in T}\frac{p_i}{P_T}\leq(Dh^{12})^{-1}.
\]
With probability at least $1-d^{-25}-e^{-cN}$,
\begin{equation}
 \sum_{i\in T}p_i\Phi(Wv_i)
 \geq cP_T\psi\!\left(\frac{\norm W_{\F}}d\right)
 \qquad\text{for every }W.
 \label{eq:tail-coercivity}
\end{equation}
\end{corollary}

\begin{proof}
Apply Lemma~\ref{lem:free-energy} to each vector $Wv_i$, then apply
Lemma~\ref{lem:weighted-query} with weights $q_i=p_i/P_T$.  Both
statements are uniform, so $W$ may depend arbitrarily on the data.
\end{proof}

\subsection{A reusable bilinear Gaussian bound}

\begin{lemma}[Bilinear Gaussian matrix bound]
\label{lem:bilinear}
Let $U,V\in\R^{d\times N}$ be independent matrices with independent
$\mathcal N(0,1/d)$ entries.  For each fixed deterministic
$M\in\R^{N\times N}$,
\begin{equation}
 \norm{UMV^\top}_{\F}\leq2\norm M_{\F}
 \label{eq:bilinear-bound}
\end{equation}
with probability at least $1-2e^{-cd}$.
\end{lemma}

\begin{proof}
Write $U=X/\sqrt d$ with $X$ standard Gaussian.  Then
\[
 \norm{UM}_{\F}^2
 =\frac1d\vecop(X)^\top(I_d\otimes MM^\top)\vecop(X).
\]
Equivalently, the quadratic-form matrix is
$d^{-1}I_d\otimes MM^\top$.  It has trace $\norm M_{\F}^2$,
Frobenius norm at most $\norm M_{\F}^2/\sqrt d$, and operator norm at
most $\norm M_{\F}^2/d$.  Lemma~\ref{lem:standard-tools}(iv), with
$x=c_0d$ and $c_0$ sufficiently small, yields
$\norm{UM}_{\F}^2\leq2\norm M_{\F}^2$ outside an event $e^{-cd}$.
Conditional on $UM$, the same argument applied to
$UMV^\top=(UM)Y^\top/\sqrt d$ gives
$\norm{UMV^\top}_{\F}^2\leq2\norm{UM}_{\F}^2$ outside another event
$e^{-cd}$.  Combining the two inequalities proves
\eqref{eq:bilinear-bound}.
\end{proof}

\begin{lemma}[Weighted centered gradient]
\label{lem:gradient}
For every deterministic vector $a=(a_1,\ldots,a_N)$,
\begin{equation}
 \norm{\sum_{i=1}^Na_ix_iv_i^\top}_{\F}
 \leq2\norm a_2
 \label{eq:gradient-bound}
\end{equation}
with probability at least $1-2e^{-cd}$.
\end{lemma}

\begin{proof}
Let $U=[u_1,\ldots,u_N]$, $V=[v_1,\ldots,v_N]$, and
\[
 M=\diag(a)-\frac1N\1 a^\top.
\]
Then $\sum_i a_ix_iv_i^\top=UMV^\top$.  The $i$th column of $M$ is
$a_i(e_i-N^{-1}\1)$, so
\[
 \norm M_{\F}^2=(1-N^{-1})\sum_i a_i^2\leq\norm a_2^2.
\]
Apply Lemma~\ref{lem:bilinear}.
\end{proof}

\subsection{The deterministic clipped comparison matrix}

Retain $b_i,R,J$ from \eqref{eq:bRJ} and define
\begin{equation}
 S_b=\sum_{i=1}^Nb_i u_iv_i^\top.
 \label{eq:Sb}
\end{equation}

\begin{lemma}[Clipped comparison]
\label{lem:comparison}
Fix $0\leq\alpha<1$ and $\beta>2$.  With probability at least
$1-d^{-20}-e^{-cN}$,
\begin{align}
 \norm{S_b}_{\F}&\leq2R,
 \label{eq:Sb-norm}\\
 \ip{\sum_i p_ix_iv_i^\top}{S_b}_{\F}&\geq\frac12J,
 \label{eq:Sb-linear}\\
 \max_{i,j}|x_j^\top S_bv_i|&\leq C,
 \label{eq:Sb-logits}\\
 \sum_i p_i\norm{S_bv_i}_2^2
 &\leq Ch\left(\sum_i p_ib_i^2+\frac{R^2}{d}\right).
 \label{eq:Sb-action}
\end{align}
Consequently, for a sufficiently small absolute $c_0>0$,
\begin{equation}
 W_0=\frac{c_0}{h}S_b
 \quad\Longrightarrow\quad
 I(W_0)\geq c\frac{J}{h}
 \geq c\frac{R^2}{Dh}.
 \label{eq:comparator-gain}
\end{equation}
\end{lemma}

\begin{proof}
The norm bound \eqref{eq:Sb-norm} follows from
Lemma~\ref{lem:bilinear} with $M=\diag(b)$.

We next prove \eqref{eq:Sb-linear}.  First remove centering and put
$G_0=\sum_i p_i u_iv_i^\top$.  Expanding gives
\begin{align*}
 \ip{G_0}{S_b}_{\F}
 &=\sum_i p_ib_i\norm{u_i}_2^2\norm{v_i}_2^2\\
 &\quad+\sum_{i<k}(p_ib_k+p_kb_i)
 (u_i^\top u_k)(v_i^\top v_k).
\end{align*}
For the diagonal term, let $a_i=p_ib_i$.  Since
$\E[\norm{u_i}_2^2\norm{v_i}_2^2]=1$ and
$\operatorname{Var}(\norm{u_i}_2^2\norm{v_i}_2^2)\leq C/d$,
\[
 \left\|\sum_i a_i
 (\norm{u_i}_2^2\norm{v_i}_2^2-1)\right\|_{L^2}
 \leq C J/\sqrt d,
\]
because $\sum_i a_i^2\leq(\sum_i a_i)^2=J^2$.  This is a centered
Gaussian polynomial of degree four, so Lemma~\ref{lem:standard-tools}(i),
with moment order $q=\lceil80\log d\rceil$, shows that it is at most
$J/8$ with failure at most $d^{-40}$.

The off-diagonal term is another centered degree-four Gaussian
polynomial.  Distinct unordered pairs are orthogonal in $L^2$, and
\[
 \E(u_i^\top u_k)^2\E(v_i^\top v_k)^2=d^{-2}.
\]
Therefore its $L^2$ norm is at most
\[
 \frac C d
 \left(\sum_{i<k}(p_ib_k+p_kb_i)^2\right)^{1/2}
 \leq C\frac{\norm p_2R}{d}.
\]
Hypercontractivity and \eqref{eq:cross-small} imply that its magnitude
is at most $J/8$ with failure at most $d^{-40}$.

Finally,
\[
 \sum_i p_ix_iv_i^\top
 =G_0-\bar u\left(\sum_i p_iv_i\right)^\top.
\]
With failure $e^{-cd}$,
\[
 \norm{\bar u}_2\leq2/\sqrt N,
 \qquad
 \norm{\sum_i p_iv_i}_2\leq2\norm p_2,
 \qquad
 \norm{S_b}_{\F}\leq2R.
\]
Hence the centering correction is at most
$8\norm p_2R/\sqrt N=o(J)$ by \eqref{eq:cross-small} and
$N\gg D$.  This proves \eqref{eq:Sb-linear}.

For \eqref{eq:Sb-logits}, let $\mathcal E_0$ denote the event
\begin{equation}
 \max_i\{\norm{u_i}_2,\norm{v_i}_2\}\leq2,
 \qquad
 \max_{i\neq j}\{|u_i^\top u_j|,|v_i^\top v_j|\}
 \leq C\sqrt{h/d}.
 \label{eq:uniform-norm-inner}
\end{equation}
Standard chi-square bounds and conditional Gaussian inner-product
tails, followed by a union bound, give
\[
 \Pp(\mathcal E_0^c)\leq d^{-40}+e^{-cd}
\]
after choosing $C$ sufficiently large.

For fixed $i,j$, define
\[
 Z_{ij}:=
 \sum_{k\notin\{i,j\}}b_k(u_j^\top u_k)(v_k^\top v_i),
 \qquad
 \mathcal A_{ij}:=
 \{\norm{u_j}_2\leq2,\ \norm{v_i}_2\leq2\}.
\]
Condition only on $(u_j,v_i)$, not on $\mathcal E_0$.  The summands in
$Z_{ij}$ are conditionally independent and centered.  Since
\[
 \norm{u_j^\top u_k}_{\psiTwo\mid u_j}
 \leq C\frac{\norm{u_j}_2}{\sqrt d},
 \qquad
 \norm{v_k^\top v_i}_{\psiTwo\mid v_i}
 \leq C\frac{\norm{v_i}_2}{\sqrt d},
\]
and the two factors are conditionally independent,
\eqref{eq:orlicz-product} shows that, on the
$\sigma(u_j,v_i)$-measurable event $\mathcal A_{ij}$,
\[
 \norm{b_k(u_j^\top u_k)(v_k^\top v_i)}
       _{\psiOne\mid u_j,v_i}
 \leq\frac{Cb_k}{d}.
\]
Consequently,
\[
 \sum_{k\notin\{i,j\}}
 \norm{b_k(u_j^\top u_k)(v_k^\top v_i)}
       _{\psiOne\mid u_j,v_i}^2
 \leq\frac{CR^2}{d^2},
 \qquad
 \max_{k\notin\{i,j\}}
 \norm{b_k(u_j^\top u_k)(v_k^\top v_i)}
       _{\psiOne\mid u_j,v_i}
 \leq\frac Cd.
\]
Conditional Bernstein with $t=50h$, followed by integration over
$(u_j,v_i)$, gives
\[
 \Pp\!\left(
 \mathcal A_{ij}\cap
 \left\{|Z_{ij}|>
 C\left(\frac{R\sqrt h}{d}+\frac hd\right)\right\}
 \right)
 \leq2e^{-50h}.
\]
Because $\mathcal E_0\subseteq\bigcap_{i,j}\mathcal A_{ij}$, a union
bound over $N^2$ pairs yields
\[
 \Pp\!\left(
 \mathcal E_0\cap
 \left\{\max_{i,j}|Z_{ij}|>
 C\left(\frac{R\sqrt h}{d}+\frac hd\right)\right\}
 \right)
 \leq2N^2e^{-50h}
 =2N^{-48}\leq d^{-40}.
\]
Moreover,
\begin{equation}
 C\left(\frac{R\sqrt h}{d}+\frac hd\right)=o(1)
 \label{eq:Sb-remainder}
\end{equation}
by \eqref{eq:R-polylog-small} and $h/d=o(1)$.  On the intersection
of $\mathcal E_0$ and the preceding remainder event, the omitted
collision terms are $O(1)$ when $i=j$ and $O(\sqrt{h/d})$ otherwise.
Thus $\max_{i,j}|u_j^\top S_bv_i|\leq C$.

Finally, on the previously established events
$\norm{\bar u}_2\leq2/\sqrt N$, $\norm{S_b}_{\F}\leq2R$, and
$\mathcal E_0$,
\[
 |\bar u^\top S_bv_i|
 \leq\norm{\bar u}_2\norm{S_b}_{\F}\norm{v_i}_2
 \leq\frac{CR}{\sqrt N}
 =C\frac Rd\frac d{\sqrt N}=o(1),
\]
where \eqref{eq:R-small} and $N\gg D$ were used.  Since
$x_j=u_j-\bar u$, this proves \eqref{eq:Sb-logits}.

To prove \eqref{eq:Sb-action}, condition on
$V=[v_1,\ldots,v_N]$ and set
\[
 B=\diag(b),
 \qquad
 P=\diag(p),
 \qquad
 K=B V^\top V P V^\top V B\succeq0.
\]
Then
\[
 \sum_i p_i\norm{S_bv_i}_2^2=\tr(UKU^\top),
 \qquad U=[u_1,\ldots,u_N].
\]
Write $U=X/\sqrt d$ and let $g\in\R^{dN}$ be the concatenation of
the $d$ rows of the standard Gaussian matrix $X$.  Then
\[
 \tr(UKU^\top)
 =g^\top B_Kg,
 \qquad
 B_K=\frac1d(I_d\otimes K)\succeq0.
\]
Since $K\succeq0$,
\[
 \tr B_K=\tr K,\qquad
 \norm{B_K}_{\F}
 =\frac{\norm K_{\F}}{\sqrt d}
 \leq\frac{\tr K}{\sqrt d},\qquad
 \norm{B_K}_{\op}
 =\frac{\norm K_{\op}}d
 \leq\frac{\tr K}d.
\]
Lemma~\ref{lem:standard-tools}(iv), with $x=c_0d$ and $c_0>0$
sufficiently small, therefore gives
\[
 \tr(UKU^\top)\leq2\tr K
\]
with conditional failure $e^{-c_0d}$.  Moreover,
\[
 \tr K=\sum_{i,k}p_ib_k^2(v_i^\top v_k)^2.
\]
On the norm event, the diagonal part is at most
$C\sum_i p_ib_i^2$.  For the off-diagonal part, Minkowski's inequality
and \eqref{eq:inner-product-moment} give, for
$q=\lceil40\log d\rceil$,
\[
 \left\|\sum_{i\neq k}p_ib_k^2(v_i^\top v_k)^2\right\|_{L^q}
 \leq\sum_{i\neq k}p_ib_k^2
 \norm{(v_i^\top v_k)^2}_{L^q}
 \leq Cq\frac{R^2}{d}.
\]
Markov's inequality proves \eqref{eq:Sb-action} with failure
$d^{-30}$.

It remains to prove the objective gain.  The centered empirical
covariance satisfies
\[
 \frac1N\sum_jx_jx_j^\top
 \preceq\frac1N\sum_ju_ju_j^\top
 \preceq\frac C d I_d
\]
with failure $e^{-cN}$ by Lemma~\ref{lem:standard-tools}(vi), taking $t=\sqrt N/4$.
Let $\eta=c_0/h$ and $W_0=\eta S_b$.  By
\eqref{eq:Sb-logits}, $|x_j^\top W_0v_i|\leq1/4$ when $c_0$ is small.
Since $N^{-1}\sum_jx_j^\top W_0v_i=0$ and
$e^z\leq1+z+Cz^2$ for $|z|\leq1/4$,
\[
 \Phi(W_0v_i)
 \leq C\eta^2\frac1N\sum_j(x_j^\top S_bv_i)^2
 \leq C\frac{\eta^2}{d}\norm{S_bv_i}_2^2.
\]
Therefore, by \eqref{eq:Sb-action},
\begin{align*}
 \sum_i p_i\Phi(W_0v_i)
 &\leq C\eta^2h
 \left(\frac1d\sum_i p_ib_i^2+\frac{R^2}{D}\right)\\
 &\leq C\eta^2hJ,
\end{align*}
where $b_i^2\leq b_i$ and \eqref{eq:J-R} were used.  On the other hand,
\eqref{eq:Sb-linear} gives
\[
 \sum_i p_ix_i^\top W_0v_i
 =\eta\ip{\sum_i p_ix_iv_i^\top}{S_b}_{\F}
 \geq\eta J/2.
\]
Since $\eta=c_0/h$, choosing $c_0$ sufficiently small makes the
free-energy term at most $\eta J/4$.  This proves
\eqref{eq:comparator-gain}.
\end{proof}

\subsection{A head frame bound and a capped-support inequality}

\begin{lemma}[Upper frame bound on the natural head]
\label{lem:head-frame}
With probability at least $1-d^{-25}$,
\begin{equation}
 \sum_{i=1}^m(x_i^\top Wv_i)^2\leq Ch\norm W_{\F}^2
 \qquad\text{for every }W\in\R^{d\times d}.
 \label{eq:head-frame}
\end{equation}
\end{lemma}

\begin{proof}
First replace $x_i$ by $u_i$.  Let
$y_i=\vecop(u_iv_i^\top)\in\R^D$.  Then
$\E[y_iy_i^\top]=I_D/D$.  Also,
\[
 \left\|\E[(y_iy_i^\top-I_D/D)^2]\right\|_{\op}\leq C/D.
\]
Indeed,
$\E[\norm{u_i}_2^2u_iu_i^\top]=((d+2)/d^2)I_d$, and similarly for
$v_i$, so
$\norm{\E[\norm{y_i}_2^2y_iy_i^\top]}_{\op}\leq C/D$.

Truncate on $\{\norm{u_i}_2,\norm{v_i}_2\leq2\}$.  The truncated
summand norm is $O(1)$, the variance proxy for the first $m\leq D$
terms is $O(m/D)\leq O(1)$, and the total truncation bias is
$O(me^{-cd})=e^{-\Omega(d)}$.  Matrix Bernstein with deviation $C_0h$
and $C_0$ sufficiently large gives
\[
 \left\|\sum_{i=1}^m y_iy_i^\top\right\|_{\op}\leq Ch
\]
with failure at most $d^{-30}$.  Hence
\[
 \sum_{i=1}^m(u_i^\top Wv_i)^2\leq Ch\norm W_{\F}^2.
\]
For the centering part,
\[
 \sum_{i=1}^m(\bar u^\top Wv_i)^2
 \leq\norm{\bar u}_2^2
 \left\|\sum_{i=1}^m v_iv_i^\top\right\|_{\op}
 \norm W_{\F}^2.
\]
Using $\norm{\bar u}_2^2\leq C/N$ and the crude bound
$\norm{\sum_{i\leq m}v_iv_i^\top}_{\op}\leq4m\leq4D$, the coefficient
is $O(D/N)=o(1)$.  Finally,
$(a-b)^2\leq2a^2+2b^2$ proves \eqref{eq:head-frame}.
\end{proof}

\begin{lemma}[Weighted capped support]
\label{lem:capped-support}
Fix $0\leq\alpha<1$.  For every $M>0$, there exists
$K=K(\alpha,M)$ such that, for every $m\geq1$ and all sufficiently
large $h$, the following holds.  If $a_1,\ldots,a_m\in\R$ satisfy
\[
 \sum_{i=1}^m a_i^2\leq B^2,
 \qquad
 B\leq\sqrt m\,h^{-K},
\]
then
\begin{equation}
 \sum_{i=1}^mp_i\min\{h,(a_i)_+\}
 \leq h^{-M}\sum_{i=1}^mp_i.
 \label{eq:capped-support}
\end{equation}
\end{lemma}

\begin{proof}
The normalization $H_{N,\alpha}^{-1}$ cancels in every ratio below.
Choose
\[
 K_0>\frac{M+3}{1-\alpha}.
\]
First suppose $m\geq2h^{K_0}$ and put
$q=\lfloor mh^{-K_0}\rfloor$, so $1\leq q<m$.  Cauchy--Schwarz gives
\begin{equation}
 \sum_{i=1}^mp_i\min\{h,(a_i)_+\}
 \leq h\sum_{i\leq q}p_i
 +B\left(\sum_{q<i\leq m}p_i^2\right)^{1/2}.
 \label{eq:cap-split}
\end{equation}
Integral comparison gives
\[
 \frac{\sum_{i\leq q}p_i}{\sum_{i\leq m}p_i}
 \leq C(q/m)^{1-\alpha}
 \leq Ch^{-K_0(1-\alpha)},
\]
so the first term in \eqref{eq:cap-split} is at most
$h^{-M-2}\sum_{i\leq m}p_i$.  Also,
\[
 \frac{(\sum_{q<i\leq m}p_i^2)^{1/2}}
 {\sum_{i\leq m}p_i}
 \leq\frac C{\sqrt m}
 \begin{cases}
 1,&0\leq\alpha<1/2,\\
 \sqrt{\log(m/q)},&\alpha=1/2,\\
 (q/m)^{1/2-\alpha},&1/2<\alpha<1.
 \end{cases}
\]
Thus the second term, divided by $\sum_{i\leq m}p_i$, is at most
\[
 Ch^{-K}\times
 \begin{cases}
 1,&\alpha<1/2,\\
 \sqrt{K_0\log h},&\alpha=1/2,\\
 h^{K_0(\alpha-1/2)},&\alpha>1/2.
 \end{cases}
\]
Choosing
$K>M+3+K_0(\alpha-1/2)_+$ proves \eqref{eq:capped-support} in this
case.

Now suppose $m<2h^{K_0}$.  Direct Cauchy--Schwarz gives
\[
 \sum_{i=1}^mp_i\min\{h,(a_i)_+\}
 \leq B\left(\sum_{i=1}^mp_i^2\right)^{1/2}.
\]
Integral comparison gives the explicit bound
\[
 \frac{\sqrt m(\sum_{i\leq m}p_i^2)^{1/2}}
 {\sum_{i\leq m}p_i}
 \leq C
 \begin{cases}
 1,&0\leq\alpha<1/2,\\
 \sqrt{\log(m+1)},&\alpha=1/2,\\
 m^{\alpha-1/2},&1/2<\alpha<1.
 \end{cases}
\]
Since $m<2h^{K_0}$, the right-hand side is at most a fixed power of
$h$.  Increasing $K$ once more proves \eqref{eq:capped-support} for the
remaining values of $m$.
\end{proof}

\section{Proof for $0\leq\alpha<1$}

\begin{proposition}[Upper bound for $\alpha<1$]
\label{prop:alpha-less-upper}
For fixed $0\leq\alpha<1$ and $\beta>2$, with probability at least
$1-d^{-c}-e^{-cN}$,
\[
 \norm{W_\star}_{\F}\leq h^C R.
\]
\end{proposition}

\begin{proof}
Work on $\mathcal E_{\rm ex}$ and on the concentration events invoked
below.  Use the deterministic diffuse tail $T$ from
Lemma~\ref{lem:strong-tail}.  Since $P_T\geq1/2$ and
$p_i/P_T\leq2/(Dh^{24})\leq(Dh^{12})^{-1}$, Corollary
\ref{cor:tail-coercivity} applies.  Put
$r=\norm{W_\star}_{\F}$ and $a=\norm{G_T}_{\F}$.  Since
$L(W_\star)\leq L(0)=h$, \eqref{eq:subset-lower} and
\eqref{eq:tail-coercivity} give
\begin{equation}
 cP_T\psi(r/d)\leq P_Hh+ar.
 \label{eq:upper-master}
\end{equation}
Lemma~\ref{lem:gradient} gives $a\leq2\sqrt{S_T}$, so
$ad\leq2h^{-A/2}$.  If $r\geq d$, then
$\psi(r/d)=r/d$, and the term $ar=(ad)(r/d)$ can be absorbed into the
left-hand side.  Hence
\[
 r/d\leq CP_Hh=o(1),
\]
contradicting $r/d\geq1$.  Therefore $r<d$ and
$\psi(r/d)=r^2/d^2$.  Solving the quadratic inequality
\[
 cr^2/d^2\leq P_Hh+ar
\]
gives
\[
 r\leq C\bigl(Da+d\sqrt{P_Hh}\bigr)
 \leq C\bigl(D\sqrt{S_T}+d\sqrt{P_Hh}\bigr).
\]
Lemma~\ref{lem:strong-tail} now yields $r\leq h^CR$.
\end{proof}

\begin{proposition}[Lower bound for $\alpha<1$]
\label{prop:alpha-less-lower}
For fixed $0\leq\alpha<1$ and $\beta>2$, with probability at least
$1-d^{-c}-e^{-cN}$,
\[
 \norm{W_\star}_{\F}\geq h^{-C}R.
\]
\end{proposition}

\begin{proof}
Work on $\mathcal E_{\rm ex}$ and on the concentration events invoked
below.  By optimality, $I(W_\star)\geq I(W_0)$ for the
deterministic-coefficient comparison matrix in
Lemma~\ref{lem:comparison}.  Hence
\begin{equation}
 I(W_\star)\geq c\frac{R^2}{Dh}.
 \label{eq:lower-gain}
\end{equation}
Let $r=\norm{W_\star}_{\F}$ and
$s_i=x_i^\top W_\star v_i$.  By \eqref{eq:deterministic-cap},
\begin{align}
 I(W_\star)
 &\leq\sum_{i\leq m}p_i\min\{h,(s_i)_+\}
 +\left|\ip{G_{\{m+1,\ldots,N\}}}{W_\star}_{\F}\right|\\
 &\leq\sum_{i\leq m}p_i\min\{h,(s_i)_+\}
 +2\sqrt{\sum_{i>m}p_i^2}\,r,
 \label{eq:lower-decomp}
\end{align}
where Lemma~\ref{lem:gradient} was used.  Since
$D^2\sum_{i>m}p_i^2\leq R^2$, the last term is at most $2Rr/D$.

Let $K_{\rm cap}=K(\alpha,3)$ be the exponent supplied by
Lemma~\ref{lem:capped-support} with $M=3$.  Fix any $K_0>2$, and then
fix
\[
 K\geq\max\{2,K_0+K_{\rm cap}+2\}.
\]
Assume for contradiction that
\begin{equation}
 r\leq c_\star R h^{-K},
 \label{eq:small-r-assumption}
\end{equation}
where $c_\star>0$ will be chosen sufficiently small.  The tail term in
\eqref{eq:lower-decomp} is at most
$2c_\star R^2/(Dh^K)$, which is at most one quarter of
\eqref{eq:lower-gain} for all sufficiently large $h$ after choosing
$c_\star$ small.

If $m=0$, this already contradicts \eqref{eq:lower-gain}.  Suppose
$m\geq1$.  Lemma~\ref{lem:head-frame} gives
\begin{equation}
 \left(\sum_{i\leq m}s_i^2\right)^{1/2}\leq C\sqrt h\,r.
 \label{eq:head-l2}
\end{equation}
If $R>\sqrt m\,h^{K_0}$, then
\eqref{eq:head-mass} and the trivial cap give
\[
 \sum_{i\leq m}p_i\min\{h,(s_i)_+\}
 \leq h\sum_{i\leq m}p_i
 \leq C\frac{hm}{D}
 \leq C\frac{R^2}{Dh^{2K_0-1}},
\]
which is at most one quarter of \eqref{eq:lower-gain} for all
sufficiently large $h$, because $K_0>2$.

If $R\leq\sqrt m\,h^{K_0}$, then \eqref{eq:head-l2} and
\eqref{eq:small-r-assumption} imply
\[
 \left(\sum_{i\leq m}s_i^2\right)^{1/2}
 \leq Cc_\star\sqrt m\,h^{K_0+1/2-K}.
\]
Our choice of $K$ implies, for all sufficiently large $h$,
\[
 Cc_\star h^{K_0+1/2-K}\leq h^{-K_{\rm cap}}.
\]
Lemma~\ref{lem:capped-support}, with $M=3$, therefore applies and yields
\[
 \sum_{i\leq m}p_i\min\{h,(s_i)_+\}
 \leq h^{-3}\sum_{i\leq m}p_i
 \leq C\frac{m}{Dh^3}
 \leq C\frac{R^2}{Dh^3},
\]
again at most one quarter of \eqref{eq:lower-gain}.  Thus the head and
tail in \eqref{eq:lower-decomp} are together strictly smaller than the
right-hand side of \eqref{eq:lower-gain}, a contradiction.  Therefore
$r\geq cRh^{-K}$.
\end{proof}

Propositions~\ref{prop:alpha-less-upper} and
\ref{prop:alpha-less-lower} prove Theorem~\ref{thm:main} for every
fixed $0\leq\alpha<1$.

\section{The endpoint $\alpha=1$}

At the endpoint,
\[
 p_i=\frac1{iH_N},
 \qquad
 H_N=\sum_{k=1}^N\frac1k
 =h+\gamma+O(N^{-1})\asymp h.
\]
Put
\[
 m_1=\left\lfloor\frac{D}{H_N}\right\rfloor.
\]
Since $D/H_N\to\infty$ and $D/(NH_N)\to0$, one has
$1\leq m_1<N$ for all sufficiently large $d$.  Moreover,
$Dp_i\geq1$ exactly when $i\leq m_1$.  Therefore
\begin{equation}
 R_F^2
 =m_1+\frac{D^2}{H_N^2}\sum_{i=m_1+1}^N\frac1{i^2}.
 \label{eq:endpoint-RF-exact}
\end{equation}
The integral bounds
\[
 \frac1{m_1+1}-\frac1{N+1}
 \leq\sum_{i=m_1+1}^N\frac1{i^2}
 \leq\frac1{m_1}
\]
and $m_1/N\to0$ imply
\[
 \sum_{i=m_1+1}^N i^{-2}\asymp m_1^{-1}.
\]
Since $m_1\asymp D/h$, \eqref{eq:endpoint-RF-exact} gives
\begin{equation}
 R_F^2\asymp\frac Dh,
 \qquad
 R_F\asymp\frac d{\sqrt h}.
 \label{eq:endpoint-scale}
\end{equation}

\subsection{A head-memory comparison}

For $K\leq N$, define
\[
 S_K=\sum_{k=1}^K u_kv_k^\top.
\]

\begin{lemma}[Head-memory event]
\label{lem:head-memory}
Let $\delta=d^{-40}$ and
$\Lambda=\log(64N^2/\delta)$.  There are absolute constants
$c,c_0,C>0$ such that, whenever
\[
 K\leq cD/h^3,
\]
then, with probability at least $1-\delta$,
\begin{align}
 (u_i-u_j)^\top S_Kv_i&\geq c_0,
 &&i\leq K,\ j\neq i,
 \label{eq:head-margin}\\
 |(u_j-u_i)^\top S_Kv_i|&\leq C/h,
 &&i>K,\ j\neq i.
 \label{eq:tail-logit}
\end{align}
\end{lemma}

\begin{proof}
Let $\mathcal E_0$ denote the event
\begin{equation}
 \max_i\bigl|\norm{u_i}_2^2-1\bigr|
 \vee\max_i\bigl|\norm{v_i}_2^2-1\bigr|\leq\frac14,
 \qquad
 \max_{i\neq j}\{|u_i^\top u_j|,|v_i^\top v_j|\}
 \leq C_0\sqrt{\Lambda/d}.
 \label{eq:endpoint-basic-event}
\end{equation}
The norm bounds follow from chi-square tails.  For a cross inner
product, condition on one vector and use the conditional Gaussian tail
bound on its norm event.  A union bound over at most $2N^2$
quantities, with $C_0$ sufficiently large, gives
\[
 \Pp(\mathcal E_0^c)\leq\delta/4.
\]

Fix $i\neq j$, let
$\mathcal F_{ij}=\sigma(u_i,u_j,v_i)$, and, for
$k\notin\{i,j\}$, define
\[
 Z_k^{(i,j)}
 =((u_i-u_j)^\top u_k)(v_k^\top v_i).
\]
Conditional on $\mathcal F_{ij}$, the variables
$\{Z_k^{(i,j)}:k\leq K,\ k\notin\{i,j\}\}$ are independent and
centered.  Their two factors are conditionally independent centered
Gaussian variables.  Hence \eqref{eq:orlicz-product} gives
\[
 \norm{Z_k^{(i,j)}}_{\psiOne\mid\mathcal F_{ij}}
 \leq C\frac{\norm{u_i-u_j}_2\norm{v_i}_2}{d}.
\]
On the $\mathcal F_{ij}$-measurable event
\[
 \mathcal A_{ij}
 :=\{\norm{u_i}_2\leq2,\ \norm{u_j}_2\leq2,\
       \norm{v_i}_2\leq2\},
\]
we consequently have
\[
 \max_{\substack{k\leq K\\k\notin\{i,j\}}}
 \norm{Z_k^{(i,j)}}_{\psiOne\mid\mathcal F_{ij}}
 \leq\frac Cd,
 \qquad
 \sum_{\substack{k\leq K\\k\notin\{i,j\}}}
 \norm{Z_k^{(i,j)}}_{\psiOne\mid\mathcal F_{ij}}^2
 \leq\frac{CK}{d^2}.
\]
Conditional Bernstein, Lemma~\ref{lem:standard-tools}(ii), with
$t=\Lambda$, implies
\begin{align*}
 &\Pp\left(
 \mathcal A_{ij}\cap
 \left\{
 \left|
 \sum_{\substack{k\leq K\\k\notin\{i,j\}}}Z_k^{(i,j)}
 \right|
 >
 C\left(\frac{\sqrt{K\Lambda}}d+\frac{\Lambda}{d}\right)
 \right\}\right)\\
 &\hspace{45mm}\leq2e^{-\Lambda}.
\end{align*}
Indeed, this follows by taking conditional probabilities given
$\mathcal F_{ij}$ and then integrating over the
$\mathcal F_{ij}$-measurable event $\mathcal A_{ij}$.  In particular,
we have not conditioned on the global event $\mathcal E_0$.

Since $\mathcal E_0\subseteq\mathcal A_{ij}$ for every pair and
$\Lambda=\log(64N^2/\delta)$, a union bound over at most $N^2$ pairs
shows that, outside an event of probability at most $\delta/32$,
simultaneously for all $i\neq j$ on $\mathcal E_0$,
\begin{equation}
 \left|
 \sum_{\substack{k\leq K\\k\notin\{i,j\}}}
 ((u_i-u_j)^\top u_k)(v_k^\top v_i)
 \right|
 \leq C\left(\frac{\sqrt{K\Lambda}}d+\frac{\Lambda}{d}\right).
 \label{eq:endpoint-remainder}
\end{equation}
Because $\delta=d^{-40}$, $N=\lceil d^\beta\rceil$, and $\beta>2$
is fixed, $\Lambda\asymp h$.  The hypothesis $K\leq cD/h^3$ gives
\[
 \frac{\sqrt{K\Lambda}}d+\frac{\Lambda}{d}
 \leq\frac{C\sqrt c}{h}+\frac{Ch}{d}
 \leq\frac Ch
\]
for all sufficiently large $d$.

Work now on the intersection of $\mathcal E_0$ and the remainder
event.  If $i\leq K$, the self term $k=i$ satisfies
\[
 (\norm{u_i}_2^2-u_j^\top u_i)\norm{v_i}_2^2
 \geq
 \left(\frac34-C_0\sqrt{\frac{\Lambda}{d}}\right)\frac34
 \geq\frac12
\]
for all sufficiently large $d$.  If $j\leq K$, the only other removed
term is
\[
 (u_i^\top u_j-\norm{u_j}_2^2)(v_j^\top v_i),
\]
whose absolute value is at most
$C\sqrt{\Lambda/d}=o(1)$.  Together with
\eqref{eq:endpoint-remainder}, this proves \eqref{eq:head-margin}, for
example with $c_0=1/4$ for all sufficiently large $d$.

If $i>K$, the self term is absent.  The only possible removed term is
the $k=j$ term when $j\leq K$, and its absolute value is at most
\[
 C\sqrt{\Lambda/d}\leq C/h,
\]
where the last inequality follows from $h^3=o(d)$.  Combining this
with \eqref{eq:endpoint-remainder} proves \eqref{eq:tail-logit}.  The
total failure probability is at most
$\delta/4+\delta/32<\delta$.
\end{proof}

Choose
\begin{equation}
 K_0=\left\lfloor cD/h^3\right\rfloor,
 \qquad
 W_{\mathrm{head}}=ThS_{K_0},
 \label{eq:head-comparator}
\end{equation}
where $T>0$ is a sufficiently large absolute constant.  On the event
of Lemma~\ref{lem:head-memory},
\begin{equation}
 L(W_{\mathrm{head}})
 \leq P_{\{K_0+1,\ldots,N\}}(h+C)+N^{-10}.
 \label{eq:head-comparator-loss}
\end{equation}
Indeed, choose $T$ so that $Tc_0\geq12$.  For $i\leq K_0$, every
wrong logit is at most $-Tc_0h$, and hence
\[
 \ell_i(W_{\mathrm{head}})
 \leq\sum_{j\neq i}e^{-Tc_0h}
 \leq N^{1-Tc_0}\leq N^{-10}.
\]
For $i>K_0$, \eqref{eq:tail-logit} gives
$\ell_i(W_{\mathrm{head}})\leq h+C$.

\subsection{Endpoint lower bound}

\begin{proposition}[Lower bound at $\alpha=1$]
\label{prop:endpoint-lower}
With probability at least $1-d^{-20}$,
\[
 \norm{W_\star}_{\F}\geq c\frac d{\sqrt h}.
\]
\end{proposition}

\begin{proof}
Work on $\mathcal E_{\rm ex}$ and on the head-memory and gradient
events specified below.  Let $C_{\rm cmp}$ be a fixed constant for
which \eqref{eq:head-comparator-loss} holds.  Choose a fixed
$s_0>C_{\rm cmp}+3$ and set
\[
 K_1=\lfloor e^{-s_0}K_0\rfloor,
 \qquad
 T_1=\{K_1+1,\ldots,N\}.
\]
Since $K_0\to\infty$, also $K_1\to\infty$.  The harmonic-number
expansion $H_n=\log n+\gamma+O(n^{-1})$ and
$K_1=e^{-s_0}K_0+O(1)$ give
\[
 H_{K_0}-H_{K_1}
 =\log(K_0/K_1)+O(K_1^{-1})
 =s_0+o(1).
\]
Moreover, $h/H_N=1+O(h^{-1})$.  Hence
\begin{equation}
 \frac h{H_N}(H_{K_0}-H_{K_1})=s_0+o(1).
 \label{eq:endpoint-harmonic-gap}
\end{equation}
Write $P_0=P_{\{K_0+1,\ldots,N\}}$.  From
\eqref{eq:head-comparator-loss} and
\eqref{eq:endpoint-harmonic-gap},
\begin{align*}
 P_{T_1}h-L(W_{\mathrm{head}})
 &\geq (P_{T_1}-P_0)h-C_{\rm cmp}P_0-N^{-10}\\
 &=\frac h{H_N}(H_{K_0}-H_{K_1})
   -C_{\rm cmp}P_0-N^{-10}\\
 &\geq s_0-C_{\rm cmp}+o(1)\geq1
\end{align*}
for all sufficiently large $d$, where $P_0\leq1$ was used.  By
optimality on $\mathcal E_{\rm ex}$,
\begin{equation}
 P_{T_1}h-L(W_\star)\geq1.
 \label{eq:endpoint-gap}
\end{equation}

The set $T_1$ is deterministic, so Lemma~\ref{lem:gradient} applies to
the deterministic coefficients $a_i=p_i\1_{\{i\in T_1\}}$.  On its
event,
\[
 \norm{G_{T_1}}_{\F}\leq2\sqrt{S_{T_1}},
 \qquad
 S_{T_1}
 =\frac1{H_N^2}\sum_{i=K_1+1}^N\frac1{i^2}
 \leq\frac C{K_1h^2}.
\]
On the other hand, \eqref{eq:subset-lower} and $\Phi\geq0$ imply
\[
 L(W_\star)
 \geq P_{T_1}h-\norm{G_{T_1}}_{\F}\norm{W_\star}_{\F}.
\]
Combining this with \eqref{eq:endpoint-gap} gives
\[
 1\leq\norm{G_{T_1}}_{\F}\norm{W_\star}_{\F}
 \leq2\sqrt{S_{T_1}}\norm{W_\star}_{\F}.
\]
Therefore
\[
 \norm{W_\star}_{\F}
 \geq\frac1{2\sqrt{S_{T_1}}}
 \geq ch\sqrt{K_1}
 \asymp\frac d{\sqrt h}.
\]
The head-memory event fails with probability at most $d^{-40}$, the
single deterministic gradient event fails with probability at most
$2e^{-cd}$, and the existence event satisfies
$\Pp(\mathcal E_{\rm ex}^c)\leq e^{-c_{\rm ex}N}$.  Their sum is at
most $d^{-20}$ for all sufficiently large $d$.
\end{proof}

\subsection{Blockwise coercivity}

\begin{lemma}[Three uniform Gaussian block estimates]
\label{lem:block-estimates}
Let
\[
 k=\lceil AD/h\rceil,
 \qquad
 B=\{k+1,\ldots,2k\},
\]
where $A=A(\beta)>0$ is a fixed sufficiently large constant.  For large $d$,
$2k\leq N$.  With probability at least $1-e^{-cd}$, all of the
following hold:
\begin{align}
 \inf_{\norm a_2=1}\max_{1\leq j\leq N}u_j^\top a
 &\geq c\sqrt{h/d},
 \label{eq:cloud-inradius}\\
 \inf_{\norm b_2=1}\frac1k\sum_{i\in B}|b^\top v_i|
 &\geq c/\sqrt d,
 \label{eq:block-l1}\\
 \norm{\sum_{i\in B}u_iv_i^\top}_{\F}
 &\leq C\sqrt k.
 \label{eq:block-bilinear}
\end{align}
\end{lemma}

\begin{proof}
We first prove \eqref{eq:cloud-inradius}.  Choose a sufficiently small
constant $\gamma=\gamma(\beta)>0$ such that
$\beta(1-2\gamma^2)>1$, and put
$r=\gamma\sqrt{h/d}$.  Let $\mathcal N$ be an $(r/8)$-net of the unit
sphere.  Then
\[
 \log|\mathcal N|\leq Cd\log d.
\]
For fixed $a\in\mathcal N$, the variables
$\sqrt d\,u_j^\top a$ are standard normal.  Mills' lower bound gives
\[
 q:=\Pp(u_j^\top a\geq2r)
 \geq c h^{-1/2}e^{-2\gamma^2h}
 =c h^{-1/2}N^{-2\gamma^2}.
\]
Therefore
\[
 \Pp\left(\max_{j\leq N}u_j^\top a<2r\right)
 \leq e^{-Nq}
 \leq\exp\!\left(-cN^{1-2\gamma^2}/\sqrt h\right).
\]
This dominates the net cardinality because
$\beta(1-2\gamma^2)>1$.  On the additional event
$\max_j\norm{u_j}_2\leq2$, every unit vector is within $r/8$ of a net
point, and the corresponding support value decreases by at most $r/4$.
Thus \eqref{eq:cloud-inradius} holds with failure $e^{-cd}$.

For \eqref{eq:block-l1}, put
$\mu_d=\E|b^\top v_i|=\sqrt{2/(\pi d)}$, which is independent of the
unit vector $b$, and define
\[
 \Delta=\sup_{\norm b_2=1}
 \left|\frac1k\sum_{i\in B}|b^\top v_i|-\mu_d\right|.
\]
Symmetrization and the contraction principle give
\[
 \E\Delta
 \leq\frac Ck\E\sup_{\norm b=1}
 \left|\sum_{i\in B}\varepsilon_i b^\top v_i\right|
 =\frac Ck\E\norm{\sum_{i\in B}\varepsilon_i v_i}_2
 \leq\frac C{\sqrt k}.
\]
Since $k\asymp d^2/h$, this is at most $\mu_d/4$ for large $d$.  As a
function of the standard Gaussian vectors $g_i=\sqrt d\,v_i$, the
quantity $\Delta$ is $(kd)^{-1/2}$-Lipschitz.  Gaussian concentration,
with $t=\mu_d/4$, yields
\[
 \Pp(\Delta>\mu_d/2)\leq e^{-ck}.
\]
This proves \eqref{eq:block-l1}.

Finally, \eqref{eq:block-bilinear} follows from
Lemma~\ref{lem:bilinear} with the diagonal coefficient matrix equal to
one on $B$ and zero elsewhere; its Frobenius norm is $\sqrt k$.
\end{proof}

\begin{lemma}[One-block coercivity]
\label{lem:block-coercivity}
For $k,B$ as in Lemma~\ref{lem:block-estimates}, with $A=A(\beta)$ sufficiently
large, with probability at least $1-e^{-cd}$,
\begin{equation}
 \frac1k\sum_{i\in B}\ell_i(W)
 \geq c\frac{\sqrt h}{d}\norm W_{\F}
 \qquad\text{for every }W\in\R^{d\times d}.
 \label{eq:block-coercivity}
\end{equation}
\end{lemma}

\begin{proof}
Work on the event of Lemma~\ref{lem:block-estimates}.  The claim is
immediate for $W=0$, since its right-hand side is zero.  Assume
$W\neq0$ and write a compact singular value decomposition
\[
 W=\sum_{a=1}^r\sigma_ax_ay_a^\top,
 \qquad \sigma_a>0.
\]
For every $i\in B$, Cauchy--Schwarz gives
\[
 \sum_{a=1}^r\sigma_a^2|y_a^\top v_i|
 \leq
 \left(\sum_{a=1}^r\sigma_a^2\right)^{1/2}
 \left(\sum_{a=1}^r\sigma_a^2(y_a^\top v_i)^2\right)^{1/2}
 =\norm W_{\F}\norm{Wv_i}_2.
\]
Sum over $i\in B$ and apply \eqref{eq:block-l1} separately to each
unit vector $y_a$.  This gives
\[
 \norm W_{\F}\frac1k\sum_{i\in B}\norm{Wv_i}_2
 \geq
 \sum_{a=1}^r\sigma_a^2
 \left(\frac1k\sum_{i\in B}|y_a^\top v_i|\right)
 \geq\frac c{\sqrt d}\norm W_{\F}^2.
\]
Division by $\norm W_{\F}>0$ yields
\begin{equation}
 \frac1k\sum_{i\in B}\norm{Wv_i}_2
 \geq c\frac{\norm W_{\F}}{\sqrt d}.
 \label{eq:average-query-norm}
\end{equation}

The cloud estimate implies, for every $y\in\R^d$,
\[
 \max_{1\leq j\leq N}u_j^\top y
 \geq c\sqrt{h/d}\,\norm y_2.
\]
For $y\neq0$, this follows by applying \eqref{eq:cloud-inradius} to
$y/\norm y_2$; for $y=0$, both sides vanish.  Consequently,
\[
 \ell_i(W)
 =\log\sum_{j=1}^N e^{(u_j-u_i)^\top Wv_i}
 \geq c\sqrt{h/d}\,\norm{Wv_i}_2-u_i^\top Wv_i.
\]
Average this inequality over $i\in B$, then use
\eqref{eq:average-query-norm}, \eqref{eq:block-bilinear}, and
Cauchy--Schwarz in Frobenius inner product:
\begin{align*}
 \frac1k\sum_{i\in B}\ell_i(W)
 &\geq c\sqrt{h/d}\,\frac1k\sum_{i\in B}\norm{Wv_i}_2
 -\frac1k\left|\ip{W}{\sum_{i\in B}u_iv_i^\top}_{\F}\right|\\
 &\geq\left(c\frac{\sqrt h}{d}-\frac C{\sqrt k}\right)\norm W_{\F}.
\end{align*}
Since $k=\lceil AD/h\rceil\geq AD/h$,
\[
 \frac C{\sqrt k}\leq\frac{C\sqrt h}{\sqrt A\,d}.
\]
Choose the fixed constant $A=A(\beta)$ sufficiently large that the
last coefficient is at most $c\sqrt h/(2d)$.  This proves
\eqref{eq:block-coercivity} simultaneously for every $W$.
\end{proof}

\begin{proposition}[Upper bound at $\alpha=1$]
\label{prop:endpoint-upper}
With probability at least $1-e^{-cd}$,
\[
 \norm{W_\star}_{\F}\leq Cdh^{3/2}.
\]
\end{proposition}

\begin{proof}
Work on $\mathcal E_{\rm ex}$ and on the block-coercivity event.  For
$i\in B$, monotonicity gives
$p_i\geq p_{2k}\asymp(kh)^{-1}$.  Lemma
\ref{lem:block-coercivity} therefore implies, for every $W$,
\[
 L(W)
 \geq p_{2k}\sum_{i\in B}\ell_i(W)
 \geq c\frac{\norm W_{\F}}{d\sqrt h}.
\]
At the minimizer, $L(W_\star)\leq L(0)=h$, and hence
$\norm{W_\star}_{\F}\leq Cdh^{3/2}$.
\end{proof}

Together with \eqref{eq:endpoint-scale}, Propositions
\ref{prop:endpoint-lower} and \ref{prop:endpoint-upper} prove
Theorem~\ref{thm:main} for $\alpha=1$.

\section{Evaluation of the clipped scale}

For $0\leq\alpha<1$, recall
$p_i\asymp N^{\alpha-1}i^{-\alpha}$.

If $0\leq\alpha<1/2$, then $\sum_i p_i^2\asymp N^{-1}$.  If clipping
occurs, its head size satisfies $m=o(D^2/N)$, so it does not change the
leading order.  Thus
\[
 R_F\asymp D/\sqrt N=d^{2-\beta/2}.
\]
At $\alpha=1/2$,
$\sum_i p_i^2\asymp(\log N)/N$; in the clipped case
$m\asymp D^2/N$ while the tail still contributes
$D^2(\log N)/N$.  Therefore
\[
 R_F\asymp d^{2-\beta/2}\sqrt{\log N}.
\]

Let $1/2<\alpha<1$.  The sharper integral asymptotic
\[
 H_{N,\alpha}=\frac{N^{1-\alpha}}{1-\alpha}(1+o(1))
\]
gives
\[
 Dp_1=(1-\alpha)d^{\,2-\beta(1-\alpha)}(1+o(1)).
\]
If $\beta(1-\alpha)>2$, then $Dp_1\to0$; if
$\beta(1-\alpha)=2$, then $Dp_1\to1-\alpha<1$.  Thus in both cases
$m=0$ for all sufficiently large $d$, and
\[
 R_F=D\norm p_2\asymp DN^{-(1-\alpha)}
 =d^{2-\beta(1-\alpha)}.
\]
If $\beta(1-\alpha)<2$, then $Dp_1\to\infty$, and the clipping
index satisfies
\[
 m\asymp(DN^{\alpha-1})^{1/\alpha}
 =d^{[2-\beta(1-\alpha)]/\alpha}.
\]
Both the head and tail in \eqref{eq:R-decomp} are of order $m$, so
\[
 R_F\asymp\sqrt m
 =d^{[2-\beta(1-\alpha)]/(2\alpha)}.
\]
The endpoint $\alpha=1$ was evaluated in
\eqref{eq:endpoint-scale}.  This proves the phase diagram in
Theorem~\ref{thm:main}.

\section{Probability ledger and logical closure}

The proof uses the following events.  They are intersected
casewise: the clipped-comparison rows are used only for $\alpha<1$,
whereas the head-memory and block rows are used only for $\alpha=1$.
\begin{center}
\begin{tabular}{lll}
\hline
Event & Content & Failure probability\\
\hline
Existence & unique finite minimizer & $e^{-c_{\rm ex}N}$\\
Key free energy & \eqref{eq:free-energy-lower}, uniformly in $y$ & $e^{-cN}+e^{-cd}$\\
Weighted query moments & \eqref{eq:weighted-cov}--\eqref{eq:weighted-fourth} & $d^{-30}$\\
Bilinear/gradient bounds & Lemmas \ref{lem:bilinear}--\ref{lem:gradient} & $e^{-cd}$\\
Clipped comparison & Lemma \ref{lem:comparison} & $d^{-20}+e^{-cN}$\\
Head frame & \eqref{eq:head-frame} & $d^{-25}$\\
Endpoint head memory & \eqref{eq:head-margin}--\eqref{eq:tail-logit} & $d^{-40}$\\
Endpoint block & \eqref{eq:block-coercivity} & $e^{-cd}$\\
\hline
\end{tabular}
\end{center}
For $\alpha<1$, take the intersection of
$\mathcal E_{\rm ex}$, the key free-energy event, the weighted-query
event for the single deterministic tail $T$, the finitely many
bilinear/gradient events used above, the clipped-comparison event, and
the head-frame event.  For $\alpha=1$, take the intersection of
$\mathcal E_{\rm ex}$, the head-memory event, the gradient event for
the deterministic set $T_1$, and the endpoint block event.  A union
bound, together with $e^{-cd}\leq d^{-c_0}$ for a suitable fixed
$c_0>0$ and all sufficiently large $d$, gives
\[
 \Pp\!\left(
 \mathcal E_{\rm ex}\cap
 \{h^{-C}R_F\leq\norm{W_\star}_{\F}\leq h^CR_F\}
 \right)
 \geq1-d^{-c_0}-e^{-cN}.
\]
Every estimate applied to $W_\star$ is deterministic or uniform over
all $W$.  The comparison matrix $W_0$ uses deterministic coefficients
$b_i$.  In the two conditional Bernstein arguments, conditioning is
only on the explicitly listed Gaussian vectors; the global norm and
inner-product events are intersected afterward.  Hence no residual
involving the adaptive minimizer is treated as independent, and the
argument does not assume the clipping behavior it proves at the
norm-scale level.

\begin{thebibliography}{9}
\bibitem{existence}
Companion note, \emph{Finite Existence and Uniqueness of the Gaussian
Softmax Associative-Memory Minimizer}, existence-and-uniqueness theorem.

\bibitem{tropp}
J.~A. Tropp,
\emph{User-Friendly Tail Bounds for Sums of Random Matrices},
Foundations of Computational Mathematics, 2012.

\bibitem{vershynin}
R.~Vershynin,
\emph{High-Dimensional Probability},
Cambridge University Press, 2018.

\bibitem{janson}
S.~Janson,
\emph{Gaussian Hilbert Spaces},
Cambridge University Press, 1997.
\end{thebibliography}

\end{document}
